\documentclass[../../main.tex]{subfiles}
\begin{document}
\section{习题}\label{section:ch06-习题}

\begin{exercise}\label{exercise:ch06-01}
用\cref{proposition:ch06-null-set-cube} 给出\cref{corollary:ch06-continuous-null} 的一个更简单的证明,不使用 Sard 定理。[提示:给定光滑映射 $F:M\to N$,构造从 $M\times\mathbb{R}^k$ 到 $N$ 的适当映射,其中 $k=\dim N-\dim M$。]
\end{exercise}

\begin{exercise}\label{exercise:ch06-02}
在特殊情形 $\partial M=\varnothing$ 下证明\cref{theorem:ch06-whitney-immersion}(Whitney 浸入定理)。[提示:不失一般性,设 $M$ 是 $\mathbb{R}^{2n+1}$ 的嵌入 $n$ 维子流形。令 $UM\subseteq T\mathbb{R}^{2n+1}$ 为 $M$ 的单位切丛(习题 5-6),并令 $G:UM\to\mathbb{RP}^{2n}$ 为映射 $G(x,v)=[v]$。用 Sard 定理推出存在某 $v\in\mathbb{R}^{2n+1}\setminus\mathbb{R}^{2n}$,使 $[v]$ 不在 $G$ 的像中,并证明从 $\mathbb{R}^{2n+1}$ 到 $\mathbb{R}^{2n}$、核为 $\mathbb{R}v$ 的投影限制为 $M$ 到 $\mathbb{R}^{2n}$ 的浸入。]
\end{exercise}

\begin{exercise}\label{exercise:ch06-03}
设 $M$ 是光滑流形, $B\subseteq M$ 是闭子集, $\delta:M\to\mathbb{R}$ 是正连续函数。证明存在光滑函数 $\widetilde{\delta}:M\to\mathbb{R}$,它在 $B$ 上为零,在 $M\setminus B$ 上为正,且处处满足 $\widetilde{\delta}(x)<\delta(x)$。[提示:考虑 $f/(f+1)$,其中 $f$ 是恰在 $B$ 上为零的光滑非负函数,并用\cref{corollary:ch06-small-positive-function}。]
\end{exercise}

\begin{exercise}\label{exercise:ch06-04}
设 $M$ 是光滑流形, $B$ 是 $M$ 的闭子集, $\delta:M\to\mathbb{R}$ 是正连续函数。
\begin{enumerate}[(a)]
\item 给定任一连续函数 $f:M\to\mathbb{R}^k$,证明存在连续函数 $\widetilde{f}:M\to\mathbb{R}^k$,它在 $M\setminus B$ 上光滑,在 $B$ 上与 $f$ 一致,且与 $f$ 是 $\delta$-接近的。[提示:用习题 6-3。]
\item 给定光滑流形 $N$ 与连续映射 $F:M\to N$,证明 $F$ 关于 $B$ 同伦于一个在 $M\setminus B$ 上光滑的映射。
\end{enumerate}
\end{exercise}

\begin{exercise}\label{exercise:ch06-05}
设 $M\subseteq\mathbb{R}^n$ 是嵌入子流形。证明 $M$ 有一个管状邻域 $U$,具有如下性质:对每个 $y\in U$, $r(y)$ 是 $M$ 中距 $y$ 最近的唯一点,其中 $r:U\to M$ 是\cref{proposition:ch06-tubular-retraction} 中定义的收缩映射。[提示:先证明若 $y\in\mathbb{R}^n$ 在 $M$ 中有最近点 $x$,则 $(y-x)\perp T_xM$。然后用\cref{theorem:ch06-tubular} 证明中的记号,证明对每个 $x\in M$ 可取 $\delta>0$,使每个 $y\in E\bigl(V_\delta(x)\bigr)$ 在 $M$ 中有最近点,且该点等于 $r(y)$。]
\end{exercise}

\begin{exercise}\label{exercise:ch06-06}
设 $M\subseteq\mathbb{R}^n$ 是紧嵌入子流形。对任一 $\varepsilon>0$,令 $M_\varepsilon$ 为 $\mathbb{R}^n$ 中到 $M$ 的距离小于 $\varepsilon$ 的点之集。证明对充分小的 $\varepsilon$, $\partial M_\varepsilon$ 是 $\mathbb{R}^n$ 中的紧嵌入超曲面,且 $\overline{M_\varepsilon}$ 是 $\mathbb{R}^n$ 中的紧正则区域,其内部含 $M$。
\end{exercise}

\begin{exercise}\label{exercise:ch06-07}
通过考虑映射 $F:\mathbb{R}\to H^2$, $F(t)=(t,|t|)$,与子集 $A=[0,\infty)\subseteq\mathbb{R}$,证明当 $M$ 具有非空边界时, \cref{theorem:ch06-approx-maps} 与\cref{corollary:ch06-extension-maps} 的结论可以不成立。
\end{exercise}

\begin{exercise}\label{exercise:ch06-08}
证明光滑流形之间每个真连续映射都同伦于真光滑映射。[提示:证明\cref{theorem:ch06-approx-maps} 证明中构造的映射 $\widetilde{F}$ 在 $F$ 为真时也是真的。]
\end{exercise}

\begin{exercise}\label{exercise:ch06-09}
令 $F:\mathbb{R}^2\to\mathbb{R}^3$ 为映射 $F(x,y)=\bigl(e^y\cos x,\ e^y\sin x,\ e^{-y}\bigr)$。对哪些正数 $r$, $F$ 横截于球面 $S_r(0)\subseteq\mathbb{R}^3$?对哪些正数 $r$, $F^{-1}\bigl(S_r(0)\bigr)$ 是 $\mathbb{R}^2$ 的嵌入子流形?
\end{exercise}

\begin{exercise}\label{exercise:ch06-10}
设 $F:N\to M$ 是横截于嵌入子流形 $X\subseteq M$ 的光滑映射,并令 $W=F^{-1}(X)$。对每个 $p\in W$,证明
\begin{align*}
T_pW=(\mathrm{d}F_p)^{-1}\bigl(T_{F(p)}X\bigr).
\end{align*}
推出:若两个嵌入子流形 $X,X'\subseteq M$ 横截相交,则对每个 $p\in X\cap X'$ 有
\begin{align*}
T_p(X\cap X')=T_pX\cap T_pX'.
\end{align*}
(用于第 146 页。)
\end{exercise}

\begin{exercise}\label{exercise:ch06-11}
设 $F:M\to N$ 与 $G:N\to P$ 是光滑映射,且 $G$ 横截于嵌入子流形 $X\subseteq P$。证明 $F$ 横截于子流形 $G^{-1}(X)$,当且仅当 $G\circ F$ 横截于 $X$。
\end{exercise}

\begin{exercise}\label{exercise:ch06-12}
设 $M$ 是紧光滑 $n$ 维流形。证明若 $N\geqslant 2n$,则每个从 $M$ 到 $\mathbb{R}^N$ 的光滑映射都可被光滑浸入一致逼近。
\end{exercise}

\begin{exercise}\label{exercise:ch06-13}
设 $M$ 是光滑流形。在本章中,我们定义了 $M$ 的两个嵌入子流形横截相交意味着什么,以及到 $M$ 中的光滑映射横截于一个嵌入子流形意味着什么。更一般地,若 $F:N\to M$ 与 $F':N'\to M$ 是到 $M$ 中的光滑映射,我们说 $F$ 与 $F'$ 互相横截,如果对每一对使 $F(x)=F'(x')$ 的 $x\in N$ 与 $x'\in N'$,空间 $\mathrm{d}F_x(T_xN)$ 与 $\mathrm{d}F'_{x'}(T_{x'}N')$ 张成 $T_{F(x)}M$。证明下列命题。
\begin{enumerate}[(a)]
\item 如上的 $N,N',F,F'$, $F$ 与 $F'$ 互相横截,当且仅当映射 $F\times F':N\times N'\to M\times M$ 横截于对角
\begin{align*}
\Delta_M=\{(x,x):x\in M\}.
\end{align*}
\item 若 $S$ 是 $M$ 的嵌入子流形,则光滑映射 $F:N\to M$ 横截于 $S$,当且仅当它横截于包含映射 $\iota:S\hookrightarrow M$。
\item 若 $F:N\to M$ 与 $F':N'\to M$ 是互相横截的光滑映射,则 $F^{-1}\bigl(F'(N')\bigr)$ 是 $N$ 的嵌入子流形,其维数等于 $\dim N+\dim N'-\dim M$。
\end{enumerate}
\end{exercise}

\begin{exercise}\label{exercise:ch06-14}
本题说明若去掉横截性假设, \cref{theorem:ch06-transverse} 可以失败到何种程度。令 $S=\mathbb{R}^n\times\{0\}\subseteq\mathbb{R}^{n+1}$,并设 $A$ 是 $S$ 的任一闭子集。证明存在真嵌入超曲面 $S'\subseteq\mathbb{R}^{n+1}$,使 $S\cap S'=A$。[提示:用\cref{theorem:ch02-closed-set-smooth}。]
\end{exercise}

\begin{exercise}\label{exercise:ch06-15}
证明\cref{theorem:ch06-global-graph}(图的整体刻画)。
\end{exercise}

\begin{exercise}\label{exercise:ch06-16}
设 $M$ 与 $N$ 是光滑流形。从 $N$ 到 $M$ 的光滑映射的一个类 $\mathcal{F}$ 称为\textbf{稳定的},如果它具有如下性质:每当 $\{F_s:s\in S\}$ 是从 $N$ 到 $M$ 的光滑映射族,且对某 $s_0\in S$ 有 $F_{s_0}\in\mathcal{F}$,则存在 $s_0$ 在 $S$ 中的邻域 $U$,使对所有 $s\in U$ 有 $F_s\in\mathcal{F}$。(粗略地说,光滑映射的一个性质若在小形变下保持,则它是稳定的。)证明若 $N$ 紧,则下列从 $N$ 到 $M$ 的光滑映射类是稳定的:
\begin{enumerate}[(a)]
\item 浸入;
\item 淹没;
\item 嵌入;
\item 微分同胚;
\item 局部微分同胚;
\item 横截于给定真嵌入子流形 $X\subseteq M$ 的映射。
\end{enumerate}
\end{exercise}

\begin{exercise}\label{exercise:ch06-17}
令 $\psi:\mathbb{R}\to\mathbb{R}$ 为紧支撑光滑函数,使 $\psi(0)=1$。用从 $\mathbb{R}$ 到 $\mathbb{R}$ 的映射族 $\{F_s:s\in\mathbb{R}\}$,其中 $F_s(x)=x\psi(sx)$,证明当 $N$ 不紧时,习题 6-16 中描述的映射类不必稳定。
\end{exercise}

\end{document}